\documentclass[11pt, notitlepage]{report}
%\usepackage[amsthm,hyperref]{ntheorem}
\usepackage[nohdr]{evan}
%\lhead{Barycentric Coordinates}
\usepackage{StoreCode}
\usepackage{titlesec}
\titleformat{\chapter}[display]
  {\normalfont\LARGE\bfseries}{\chaptertitlename\ \thechapter}{10pt}{\huge}
\titlespacing*{\chapter}{0pt}{0pt}{10pt}
%\renewcommand{\vec}[1]{\overrightarrow{#1}}

\newif\iflong
\longtrue


%TODO LIST:
%Add special circles
%Clean up TSTST #2

\begin{document}
\title{Barycentric Coordinates in Olympiad Geometry}
\author{
  Max Schindler\thanks{Mewto55555, Missouri.
  I can be contacted at \href{mailto:igoroogenflagenstein@gmail.com}{igoroogenflagenstein@gmail.com}.}
  \qquad
  Evan Chen\thanks{v\_Enhance, SFBA.
  %I am much indebted to Max Schindler for this fantastic technique.
  I can be reached at \href{mailto:chen.evan6@gmail.com}{chen.evan6@gmail.com}.}
  }
\date{July 13, 2012}

\addcontentsline{toc}{chapter}{Title Page}
\maketitle


%\vspace*{.2\textheight}
\begin{center}
\begin{minipage}{8.4cm}
  \textit{I suppose it is tempting, if the only tool you have is a hammer, to treat everything as if it were a nail.}
\end{minipage}
\end{center}

%\vspace*{2cm}

\begin{abstract}
In this paper we present a powerful computational approach to large class of olympiad geometry problems-- barycentric coordinates.  We then extend this method using some of the techniques from vector computations to greatly extend the scope of this technique.

Special thanks to Amir Hossein and the other olympiad moderators for helping to get this article featured: I certainly did not have such ambitious goals in mind when I first wrote this!
\end{abstract}
%\clearpage
%\addcontentsline{toc}{chapter}{Table of Contents}
\tableofcontents

\chapter{Preliminaries}
\begin{quote}
  So many problems are killed by it.  It's just that no one knows it.
  \\ --Max Schindler
\end{quote}

Over the course of olympiad geometry, several computational approaches have surfaced as a method of producing complete solutions to geometry problems given sufficient computational fortitude.  Each has their advantages and drawbacks.

\textit{Barycentric coordinates}, also called areal coordinates, provide a new ``bash'' approach for geometry problems.  Barycentric coordinates offer a length-based, coordinate approach to geometry problems.

\section{Advantages of barycentric coordinates}
The advantages of the system include
\begin{itemize}
  \item Sides of the triangle playing the role of the axes.
  \item Simple expressions for lines in general, making it computationally feasible to intersect lines.
  \item Simple forms for some common points (centroid, incenter, symmedian point...)
  \item Very strong handling of ratios of lengths.
  \item A useful method for dropping arbitrary perpendiculars.
  \item An area formula.
  \item Circle formula.
  \item Distance formula.
\end{itemize}


This arsenal of tools is far more extensive than that of many other computational techniques.  Cartesian coordinates are woefully inadequate for most olympiad geometry problems because the forms for special points are typically hideous, and the equation of a circle is difficult to work with.  Complex numbers and vectors are more popular, but the concept of an ``equation of a line'' is complicated in the former and virtually nonexistent in the latter.


\section{Notations and Conventions}
\begin{StoreCode}{tech}
Throughout this paper, $\triangle ABC$ is a triangle with vertices in counterclockwise order.   The lengths will be abbreviated $a=BC$, $b=CA$, $c=AB$.  These correspond with points in the vector plane $\vec A$, $\vec B$, $\vec C$.
%Finally, we will assume that the circumcenter $O$ of $ABC$ is at the null vector $\vec 0$. %Do we want this?

For arbitrary points $P$, $Q$, $R$, $[PQR]$ will denote the signed area of $\triangle PQR$.\footnote{For $ABC$ counterclockwise, this is positive when $P$, $Q$ and $R$ are in counterclockwise order, and negative otherwise.  When $ABC$ is labeled clockwise the convention is reversed; that is, $[PQR]$ is positive if and only if it is oriented in the same way as $ABC$.  In this article, $ABC$ will always be labeled counterclockwise.}
\end{StoreCode}


\section{How to Use this Article}
You could read the entire thing, but the page count makes this prospect rather un-inviting.  About half of it is theory, and half of it is example problems; it's probably possible to read either half only and then go straight to the problems.

I'm assuming that the majority of readers would like to read just the examples, and attack the problems with a formula sheet beside them.  I've included Appendix \ref{appendix:formulasheet} to facilitate these needs; I also have a version which includes \textit{just} the formula sheet and the problems, which should be floating around.

Happy bashing!

\chapter{The Basics}
\begin{quote}
  We will bary you!
  \\ -- Schindler-Khruschev
\end{quote}

\section{The Coordinates}
\begin{StoreCode}{tech}
\begin{definition}
  Each point in the plane is assigned an ordered triple of real numbers $P = (x,y,z)$ such that \[ \vec P = x \vec A + y \vec B + z \vec C \quad \text{and} \quad x+y+z = 1 \]
\end{definition}
\end{StoreCode}

It is not hard to verify that the coordinates of any point are well-defined.  These are sometimes called areal coordinates because if $P=(x,y,z)$, then the signed area $[CPB]$ is equal to $x[ABC]$, and so on \cite{infinity,abel}.  In other words, these coordinates can viewed as \[ P = \frac{1}{[ABC]}([PBC], [PCA], [PAB]) \]

Of course, notice that $A=(1,0,0)$, $B=(0,1,0)$ and $C=(0,0,1)$!  This is why barycentric coordinates are substantially more suited for standard triangle geometry problems.

\begin{exercise}
  Find the coordinates for the midpoint of $AB$.
\end{exercise}
\begin{exercise}
  When is a point $(x,y,z)$ in the interior of a triangle?
\end{exercise}


\section{Lines}
\subsection{The Equation of a Line}
The equation of a line \cite{infinity,abel} in barycentric coordinates is astoundingly simple.
\begin{StoreCode}{tech}
\begin{theorem}[Line]
  \label{tech:line}
  The equation of a line is $ux+vy+wz = 0$ where $u,v,w$ are reals.  (These $u$, $v$ and $w$ are unique up to scaling.)
\end{theorem}
\end{StoreCode}
This is a corollary of the area formula, Theorem \ref{tech:area}.

In particular, if a line $\ell$ passes through a vertex, say $A$, then $u(1) + v(0) + w(0) = 0 \implies u= 0$.  So we rearrange to obtain
\begin{corollary}[Line through a vertex]
  The equation of a line passing through $A$ is simply of the form $y = kz$ for some constant $k$.
\end{corollary}

In particular, the equation for the line $AB$ is simply $z=0$, by substituting $(1,0,0)$ and $(0,1,0)$ into $ux+vy+wz=0$.

\subsection{Ceva and Menelaus}
In fact, the above techniques are already sufficient to prove both Ceva's and Menelaus's Theorem, as in \cite{infinity}.
\begin{corollary}[Ceva's Theorem]
  \label{tech:ceva}
  Let $AD$, $BE$ and $CF$ be cevians of a triangle $ABC$.  Then the cevians concur if and only if \[ \frac{BD}{DC} \frac{CE}{EA} \frac{AF}{FB} = 1\]
\end{corollary}
\begin{proof}
  Since $D$ lies on $BC$, the point $D$ has the form $D=(0,d,1-d)$.  So the equation of line $AD$  is simply \[ z = \frac{1-d}{d} y \] Similarly, if we let $E=(1-e,0,e)$ and $F=(f,1-f,0)$ then the lines $BE$ and $CF$ have equations $x = \frac{1-e}{e} z$ and $y = \frac{1-f}{f} x$, respectively.

  Notice that this system of three equations is homogeneous, so we may ignore the condition that $x+y+z=1$ temporarily.  Then it is easy to see that this equation has solutions if and only if \[ \frac{(1-d)(1-e)(1-f)}{def} = 1 \] which is equivalent to Ceva's Theorem.
\end{proof}

\begin{exercise}
  Prove Menelaus's Theorem.
\end{exercise}

\section{Special points in barycentric coordinates}
\label{sec:specialpt}
Here we give explicit forms for several special points in barycentric coordinates (compiled from \cite{infinity,wolfram,yiu}).  In this section, it will be understood that $(u:v:w)$ refers to the point $\frac{1}{u+v+w} (u,v,w)$; that is, we are not normalizing the coordinates such that they sum to $1$.\footnote{These ``standard'' coordinates with $u+v+w=1$ are sometimes referred to \textit{homogenized barycentric coordinates} to emphasize the distinction.}

\StoreText{\subsubsection*{Coordinates of Special Points} From this point on, the point $(kx:ky:kz)$ will refer to the point $(x,y,z)$ for $k \neq 0$.  In fact, the equations for the line and circle are still valid; hence, when one is simply intersecting lines and circles, it is permissible to use these un-homogenized forms in place of their normal forms.}{tech}

\begin{StoreCode}{tech}
Again, the coordinates here are not homogenized!

\begin{tabular}{r|ll}
Point & Coordinates & Sketch of Proof \\ \hline
Centroid & $G = (1:1:1)$ & Trivial \\
Incenter & $I = (a:b:c)$ & Angle bisector theorem \\
Symmedian point & $K=(a^2:b^2:c^2)$ & Similar to above \\
Excenter & $I_a = (-a:b:c)$, etc. & Similar to above \\
Orthocenter & $H=(\tan A : \tan B : \tan C)$ & Use area definition \\
Circumcenter & $O=(\sin 2A : \sin 2B : \sin 2C)$ & Use area definition
\end{tabular}

\end{StoreCode}

One will notice that $O$ and $H$ are not particularly nice in barycentric coordinates (as compared to in, say, complex numbers), but $I$ and $K$ are particularly pretty.

\begin{StoreCode}{tech}
  If absolutely necessary, it is sometimes useful to convert the trigonometric forms of $H$ and $O$ into expressions entirely in terms of the side lengths (cf. \cite{wolfram,yiu}) by \[ O = (a^2(b^2+c^2-a^2) : b^2(c^2+a^2-b^2) : c^2(a^2+b^2-c^2)) \] and \[ H = ((a^2+b^2-c^2)(c^2+a^2-b^2) : (b^2+c^2-a^2)(a^2+b^2-c^2) : (c^2+a^2-b^2)(b^2+c^2-a^2)) \]
\end{StoreCode}
%\StoreText{\eject}{tech}
As we will see in section \ref{sec:conway}, Conway's notation gives a more viable way of manipulating the expressions.

\begin{exercise}
  Prove the formulas given in the table above are correct (or email me telling me about a typo.)
\end{exercise}

Some of these proofs are given in \cite{abel}.

When calculating ratios and areas, etc., it is important that coordinates are normalized.  There are circumstances where it is both legitimate and computationally easier to use $(kx:ky:kz)$ in place of $(x,y,z)$.  See section \ref{sec:notnormal} for details.  It is recommended that one keep coordinates normalized when they are still learning this technique in order to avoid pitfalls.

\chapter{Standard Strategies}
\begin{quote}
  Uh, perpendiculars should be okay right?  Dot product?
  \\ -- Evan Chen, on EFFT
\end{quote}

This chapter will cover some typical techniques which arise in olympiad problems.

First, we begin the following definition \cite{infinity}:
\begin{StoreCode}{tech}
\begin{definition}
  The \textit{displacement vector} of two (normalized) points $P = (p_1, p_2, p_3)$ and $Q = (q_1, q_2, q_3)$ is denoted by $\overrightarrow{PQ}$ and is equal to $(p_1-q_1, p_2-q_2, p_3-q_3)$.
\end{definition}
\end{StoreCode}

Note that the sum of the coordinates of a displacement vector is $0$.

In the proofs of this chapter, we will invoke some properties which will be familiar to those who have done vector bashing; namely, when $\vec O = \vec 0$
\begin{align*}
  \vec A \cdot \vec A &= R^2 \\
  \vec A \cdot \vec B &= R^2-c^2/2
\end{align*}
where $R$ is the circumradius.

If you don't care about the proofs, then this doesn't matter; otherwise, see appendix \ref{appendix:vectors} for a treatise on some elementary vector bashing.

This chapter will frequently involve translating $O$ to the null vector $\vec 0$, so that we may invoke useful properties about dot products mentioned above.  This is valid since the point $(x,y,z)$ satisfies $x+y+z=1$, so translation is actually an invariant.  As a consequence, however, take heed of the following.
\begin{note}
  It is important that $x+y+z=1$ when doing calculations!
\end{note}





\StoreText{\subsubsection*{A Note on Scaling Displacement Vectors} In EFFT, one can write a displacement vector $(x,y,z)$ as $(kx:ky:kz)$, and the theorem will still be true.  This is also true for Strong EFFT, but NOT for the distance formula.}{tech}

\section{EFFT: Perpendicular Lines}
In fact there is a fast characterization for when two lines $MN$ and $PQ$ are perpendicular.  This lemma also appears in \cite{book}, with a proof via the Pythagorean Theorem.
\begin{StoreCode}{tech}
\begin{theorem}[Evan's Favorite Forgotten Trick]
  \label{tech:efft}
  Consider displacement vectors $\overrightarrow{MN} = (x_1, y_1, z_1)$ and $\overrightarrow{PQ} = (x_2, y_2, z_2)$.  Then $MN \perp PQ$ if and only if
  \[  0 = a^2(z_1y_2 + y_1z_2) + b^2(x_1z_2 + z_1x_2) + c^2(y_1x_2 + x_1y_2)\]
\end{theorem}
\end{StoreCode}
\begin{proof}
  Translate $\vec O$ to $\vec 0$.  It is necessary and sufficient that $(x_1\vec A + y_1\vec B + z_1 \vec C) \cdot (x_2 \vec A + y_2 \vec B + z_2 \vec C) = 0$.  Expanding, this is just
  \[ \cycsum \left(x_1x_2 \vec A \cdot \vec A \right) + \cycsum \left( \left(x_1y_2 + x_2y_1\right) \vec A \cdot \vec B \right) = 0 \]

  Utilizing the fact that $\vec O = \vec 0$, we make the standard substitutions and rearrange:

  \begin{align*}
    \iff \cycsum (x_1x_2 R^2) + \cycsum \left(x_1y_2 + x_2y_1\right)\left(R^2 - \frac{c^2}{2}\right) &= 0 \\
    \iff R^2 \left( \cycsum (x_1x_2) + \cycsum(x_1y_2 + x_2y_1) \right) &= \half \cycsum \left( (x_1y_2 + x_2y_1)(c^2) \right) \\
    \iff R^2 (x_1 + y_1 + z_1)(x_2 + y_2 + z_2) &= \half \cycsum \left( (x_1y_2 + x_2y_1)(c^2) \right) \\
    \iff R^2 \cdot 0 \cdot 0 &= \half \cycsum \left( (x_1y_2 + x_2y_1)(c^2) \right) \\
    \iff 0 &= \cycsum \left( (x_1y_2 + x_2y_1)(c^2) \right)
  \end{align*}
  since $x_1+y_1+z_1 = x_2+y_2+z_2 = 0$ (remember that $(x_1, y_1, z_1)$ is a displacement vector!).
\end{proof}
We will see later, in Section \ref{sec:notnormal}, that this holds for \textit{scaled} displacement vectors.

This will be abbreviated EFFT\footnote{Hi Simon.} throughout this document.  As we shall see, EFFT increases the scope of solvable problems by providing a usable method for constructing perpendiculars, and in particular, perpendicular bisectors.

EFFT is particularly useful when one of the displacement vectors is the side of the triangle.  For instance, we may deduce
\begin{StoreCode}{tech}
\begin{corollary}
  Consider a displacement vector $\overrightarrow{PQ} = (x_1,y_1,z_1)$.  Then $PQ \perp BC$ if and only if
  \[ 0 = a^2(z_1-y_1) + x_1(c^2-b^2) \]
\end{corollary}
\begin{corollary}
  \label{tech:perpbis}
  The perpendicular bisector of $BC$ has equation
  \[ 0 = a^2(z-y) + x(c^2-b^2) \]
\end{corollary}
\end{StoreCode}
\begin{exercise}
  Prove the above corollaries, and deduce the Pythagorean Theorem.
\end{exercise}

EFFT is used in many problems, e.g. problems \ref{exam:2008usamo2} and \ref{exam:2012wootpo47}.  Its generalization Strong EFFT (cf. Theorem \ref{tech:strongefft}) is also useful and is invoked in problem \ref{exam:2005islg5}.






\section{Distance Formula}
The distance formula \cite{infinity} can be proved by precisely the same means as EFFT:
\begin{StoreCode}{tech}
\begin{theorem}[Distance Formula]
  \label{tech:distance}
  Consider a displacement vector $\overrightarrow{PQ} = (x,y,z)$.  Then \[ |PQ|^2 = -a^2yz - b^2zx - c^2xy \]
\end{theorem}
\end{StoreCode}
\begin{exercise}
  Prove the distance formula by using \[ PQ^2 = (x\vec A + y\vec B + z \vec C) \cdot (x \vec A + y \vec B + z \vec C) \] and the fact that $x+y+z=0$.
\end{exercise}









\section{Circles}

\subsection{Equation of the Circle}
Not surprisingly, the equation of a circle \cite{infinity} is much more annoying than anything that has showed up so far.  Nonetheless, it is certainly usable, particularly in certain situations described below.

\begin{StoreCode}{tech}
\begin{theorem}
  \label{tech:circle}
  The general equation of a circle is \[ -a^2yz - b^2zx -c^2xy + (ux+vy+wz)(x+y+z) = 0 \] for reals $u,v,w$.
\end{theorem}
\end{StoreCode}
\begin{proof}
  Assume the circle has center $(j,k,l)$ and radius $r$.  Then this is just
  \[ -a^2(y-k)(z-l) - b^2(z-l)(x-j) - c^2(x-j)(y-k) = r^2 \]
  Expand everything, and collect terms to get
  \[ -a^2yz - b^2zx - c^2xy + C_1 x + C_2 y + C_3 z = C \]
  for some hideous constants $C_i$ and $C$.  Since $x+y+z = 1$, we can rewrite this as
  \[ -a^2yz - b^2zx - c^2xy + ux + vy + wz = 0 \iff -a^2xy - b^2yz -c^2zx + (ux+vy+wz)(x+y+z) = 0\]
  where $u=C_1-C$, et cetera.
\end{proof}

Fortunately, most circles that we are concerned with have much cleaner forms.  For example, as in \cite{infinity},
\begin{StoreCode}{tech}
\begin{corollary}
  The circumcircle has equation \[ a^2yz + b^2zx + c^2xy = 0 \]
\end{corollary}
\end{StoreCode}
\begin{exercise}
  Prove the above corollary.
\end{exercise}

Furthermore, circles passing through vertices of the triangle make for more elegant circle formulas, since many terms in the circle equation vanish.

\begin{note}
  Notice that the equation of the circle is homogeneous, so it is possible to scale $(x,y,z)$ arbitrarily and still have a true equation.  See section \ref{sec:notnormal}.
\end{note}

\chapter{Trickier Tactics}
\begin{quote}
  All you have to do is construct a parallelogram!
  \\ -- James Tao
\end{quote}

Here are some strategies which are occasionally useful.
\section{Areas and Lines}
The area formula \cite{infinity,abel,yiu} is as follows:

\begin{StoreCode}{tech}
\begin{theorem}[Area Formula]
  \label{tech:area}
  The area of a triangle with vertices $P=(x_1,y_1,z_1)$, $Q=(x_2,y_2,z_2)$ and $R=(x_3,y_3,z_3)$ is
  \[ [PQR] = [ABC] \cdot \left| \begin{array}{ccc} x_1 & y_1 & z_1 \\ x_2 & y_2 & z_2 \\ x_3 & y_3 & z_3 \end{array} \right| \]
\end{theorem}
\end{StoreCode}

A proof of the area formula can be found in \cite{abel}.

From this, we immediately find two corollaries, which also appear in \cite{infinity} and \cite{abel}.
\begin{StoreCode}{tech}
\begin{corollary}[First Collinearity Criteria]
  \label{tech:collnaive}
  The points $P=(x_1:y_1:z_1)$, $Q=(x_2:y_2:z_2)$ and $R=(x_3:y_3:z_3)$ are collinear if and only if
  \[ \left| \begin{array}{ccc} x_1 & y_1 & z_1 \\ x_2 & y_2 & z_2 \\ x_3 & y_3 & z_3 \end{array} \right| = 0 \]
\end{corollary}
\end{StoreCode}
\begin{proof}
  The points are collinear iff the area of the triangle they determine is zero, and this doesn't change under scaling $(x,y,z) \mapsto (kx:ky:kz)$.
\end{proof}

\begin{StoreCode}{tech}
\begin{corollary}[Line Through 2 Points]
  \label{tech:line2pt}
  The equation of a line through the points $P=(x_1:y_1:z_1)$ and $Q=(x_2:y_2:z_2)$ is
  \[ \left| \begin{array}{ccc} x_1 & y_1 & z_1 \\ x_2 & y_2 & z_2 \\ x & y & z \end{array} \right| = 0\]
\end{corollary}
\end{StoreCode}
The previous corollary also implies the line formula.

On the other hand, when normalized, Corollary \ref{tech:collnaive} can be simplified by
\begin{StoreCode}{tech}
\begin{corollary}[Second Collinearity Criteria]
  \label{tech:coll}
  The points $P=(x_1,y_1,z_1)$, $Q=(x_2,y_2,z_2)$ and $R=(x_3,y_3,z_3)$, %where $x_i+y_i+z_i=1$ for all $i$,
  are collinear if and only if
  \[ \left| \begin{array}{ccc} x_1 & y_1 & 1 \\ x_2 & y_2 & 1 \\ x_3 & y_3 & 1 \end{array} \right| = 0\]
  Cyclic variations hold.
\end{corollary}
\end{StoreCode}
\begin{proof}
  It is a property of the determinant that
  \[ \left| \begin{array}{ccc} x_1 & y_1 & z_1 \\ x_2 & y_2 & z_2 \\ x_3 & y_3 & z_3 \end{array} \right| = \left| \begin{array}{ccc} x_1 & y_1 & x_1+y_1+z_1 \\ x_2 & y_2 & x_2+y_2+z_2 \\ x_3 & y_3 & x_3+y_3+z_3 \end{array} \right|\]
\end{proof}

This has a substantial computational advantage.




\section{Non-normalized Coordinates}
\label{sec:notnormal}

Suppose $x+y+z=1$.  Non-normalized coordinates involve the notion of extending the notation $(x,y,z)$ as follows: for all nonzero $k$, the triple $(kx:ky:kz)$ refers to the point $(x,y,z)$.  We will also make the same scaling for displacement vectors: the displacement vector $(x,y,z)$ can also be written $(kx:ky:kz)$.

It turns out that, if one is merely intersecting lines and circles, we can actually use non-normalized coordinates during the computation, simplifying calculations greatly (cf. \cite{infinity}).  This holds because the equations for lines and circles are \textit{homogeneous}; that is, scaling $x,y,z$ in them does not affect the validity of the equation.  However, when using the area formula, distance formula, etc. it is extremely important that points be normalized before the calculations are applied.

\begin{exercise}
  Check that the distance formula fails for non-normalized coordinates, but holds for EFFT.
\end{exercise}

%From here on outwards, the point $(kx:ky:kz)$ will always refer to the point $(x,y,z)$ unless otherwise specified.
Some sources, including the first version of this article, simply use $(kx,ky,kz)$ instead of $(kx:ky:kz)$.  We use the latter to avoid confusion, but a reader should keep in mind that not all authors choose to do this.






\section{O, H, and Strong EFFT}
As we can see easily in Section \ref{sec:specialpt}, in fact the least elegant points are $H$ and $O$.  This seems contrary to expectation: in complex numbers and vectors, we frequently set $\vec O$ as zero, and $\vec H = \vec A + \vec B + \vec C$.  One might be inclined to think that, given the vector-based definition of barycentric coordinates, we might be able to utilize this.

Unfortunately, the area formula, and hence the line equation, relies on the fact that $(x,y,z)$ reflect the actual ratio of areas.  However, we can get some partial tools to simplify computations related to $O$ and $H$.

\begin{StoreCode}{tech}
\label{tech:strongefft}
\begin{theorem}[Strong EFFT]
  Suppose $M$, $N$, $P$ and $Q$ are points with
  \begin{align*}
    \overrightarrow{MN} &= x_1 \overrightarrow{AO} + y_1 \overrightarrow{BO} + z_1 \overrightarrow{CO} \\
    \overrightarrow{PQ} &= x_2 \overrightarrow{AO} + y_2 \overrightarrow{BO} + z_2 \overrightarrow{CO}
  \end{align*}
  If either $x_1+y_1+z_1=0$ or $x_2+y_2+z_2=0$, then $MN \perp PQ$ if and only if
  \[  0 = a^2(z_1y_2 + y_1z_2) + b^2(x_1z_2 + z_1x_2) + c^2(y_1x_2 + x_1y_2)\]
\end{theorem}
\end{StoreCode}
\begin{proof}
  Reproduce the proof of Theorem \ref{tech:efft}, with a few modifications.  Here, the ``displacement vector'' $\vec PQ$ has already been shifted, so that the properties of the dot product $A \cdot A = R^2$ and $A \cdot B = R^2 - c^2/2$ hold.

  In the final step, instead of $0 \cdot 0 = 0$ we have $0 \cdot (x_2+y_2+z_2) = 0$, which is just as well.
\end{proof}
\begin{note}
  If $MN = (x_1,y_1,z_1)$ is a displacement vector, then we automatically have both $x_1+y_1+z_1=0$ and $\vec MN = x_1 \vec A + y_1 \vec B + z_1 \vec C$, since this is independent of the location of the null vector (here $\vec O$).  Therefore, Strong EFFT is usually useful for showing some more unusual vector is perpendicular to a standard displacement vector.

  Once again, scaling is valid: the condition is homogeneous in each of $x_1,y_1,z_1$ and $x_2,y_2,z_2$.
\end{note}

This is particularly useful if one of the points involved is $H$, for $\vec H = \vec A + \vec B + \vec C$.  See problem \ref{exam:2005islg5} for an example.

The other point for which Strong EFFT is useful is $O$; for example, the following corollary is immediate
\begin{StoreCode}{tech}
\begin{corollary}
  The equation for the tangent to the circumcircle at $A$ is $b^2z+c^2y=0$.
\end{corollary}
\end{StoreCode}
\begin{exercise}
  Prove the above corollary.
\end{exercise}







\section{Conway's Formula}
\label{sec:conway}
Conway's Formula gives a nice way to work with (some) angles in a triangle.  Most of the results here are from \cite{spanish}.

We start by defining some notation:

\begin{StoreCode}{tech}
\begin{definition}[Conway's Notation]
  \label{tech:conwaydef}
  Let $S$ be twice the area of the triangle.  Define $S_{\theta}=S\cot{\theta}$, and define the shorthand notation $S_{\theta\phi}=S_{\theta}S_{\phi}$.
\end{definition}
\end{StoreCode}

\begin{StoreCode}{tech}
\begin{fact}
We have $S_A=\frac{-a^2+b^2+c^2}{2}=bc\cos{A}$ and its cyclic variations.  (We also have $S_{\omega}=\frac{a^2+b^2+c^2}{2}$, where $\omega$ is the Brocard angle.  This follows from $\cot{\omega}=\cot{A}+\cot{B}+\cot{C}$.)
\end{fact}
\end{StoreCode}

\begin{StoreCode}{tech}
\begin{fact}
  We have the identities \[S_B + S_C = a^2 \] and \[ S_{AB}+S_{BC}+S_{CA}=S^2 \]
\end{fact}
\end{StoreCode}
\begin{exercise}
  Prove the above facts.
\end{exercise}

This gives us slightly more convenient ways to express some things.  For example

\begin{StoreCode}{tech}
\begin{fact}
  $O=(a^2S_A : b^2S_B : c^2S_C)$ and $H=(S_{BC} : S_{CA} : S_{AB}) = \left( \frac{1}{S_A} : \frac{1}{S_B} : \frac{1}{S_C} \right)$.
\end{fact}
\end{StoreCode}

(Note that these do not sum to $1$; we are using the conventions spelled out in section \ref{sec:notnormal}).

Now for the main result of this chapter:
\begin{StoreCode}{tech}
\begin{theorem}[Conway's Formula]
  \label{tech:conwayformula}
   Given a point $P$ with counter-clockwise directed angles $\dang PBC=\theta$ and $\dang BCP=\phi$, we have $P=(-a^2 : S_C+S_\phi : S_B+S_\theta)$.
\end{theorem}
\end{StoreCode}

\begin{proof}
  By the Law of Sines on triangle $PBC$, $BP=\frac{a\sin{\phi}}{\sin{(\theta+\phi)}}$, and $CP=\frac{a\sin{\theta}}{\sin{(\theta+\phi)}}$, so the area of triangle $PBC$ is \[ \frac{1}{2} \cdot BC \cdot BP \sin{\theta} = \frac{a^2\sin{\theta}\sin{\phi}}{2\sin{\left(\theta+\phi\right)}} \]

  Similarly, the area of $PCA$ and $PBA$ can be found to be $\frac{ab\sin{\theta}\sin{\left(\phi+C\right)}}{2\sin{(\theta+\phi)}}$ and $\frac{ac\sin{\phi}\sin{\left(\theta+B\right)}}{2\sin{(\theta+\phi)}}$ respectively.

  Eliminating the denominator gives that $P$'s barycentric coordinates are at:
  \begin{align*}
    & (-a^2\sin{\theta}\sin{\phi} : ab\sin{\theta}\sin{\left(\phi+C\right)} : ac\sin{\phi}\sin{\left(\theta+B\right)}) \\
    &=\left(-a^2 : ab\sin{\left(\phi+C\right)}{\sin{\phi}} : \frac{ac\sin{(\theta+B)}}{\sin{\theta}}\right) \\
    &=(-a^2 : ab\cos{C}+ab\sin{C}\cot{\phi} : ca\cos{B}+ca\sin{B}\cot{\theta}) \\
    &=(-a^2 : S_C+S_\phi : S_B+S_\theta)
  \end{align*}
  as desired.
\end{proof}

An example of an application of Conway's Formula can be found in problem \ref{exam:2001islg1}.

\section{A Few Final Lemmas}
\begin{StoreCode}{tech}
\begin{lemma}[Parallelogram Lemma]
  The points $ABCD$ form a parallelogram iff $A+C = B+D$ (here the points are normalized), where addition is done component-wise.
\end{lemma}
\begin{lemma}[Concurrence Lemma]
  The three lines $u_ix+v_iy+w_iz=0$, for $i=1,2,3$ are concurrent if and only if
  \[ \left| \begin{array}{ccc} u_1 & v_1 & w_1 \\ u_2 & v_2 & w_2 \\ u_3 & v_3 & w_3 \end{array} \right| = 0\]
\end{lemma}
\end{StoreCode}

\chapter{Example Problems}
\begin{quote}
  Graders received some elegant solutions, some not-so-elegant solutions, and some so not elegant solutions.
  \\ -- Delong Meng
\end{quote}

\newcommand{\secps}{\subsection{Problem and Solution}}
\newcommand{\seccm}{\subsection{Commentary}}

Let us now demonstrate the power of the strategy on a few... victims.


\newcommand{\barydiagram}[2]{
\begin{figure}[ht]%
\centering%
\includegraphics{diagram/#2.pdf}%
\caption{#1}%
\label{fig:#1}%
\end{figure}
}
\section{USAMO 2001/2}
\label{exam:2001usamo2}
We begin with a USAMO problem which literally (figuratively?) screams for an analytic solution.

\secps
\begin{problem}[USAMO 2001/2]
  Let $ABC$ be a triangle and let $\omega$ be its incircle. Denote by $D_1$ and $E_1$ the points where $\omega$ is tangent to sides $BC$ and $AC$, respectively. Denote by $D_2$ and $E_2$ the points on sides $BC$ and $AC$, respectively, such that $CD_2=BD_1$ and $CE_2=AE_1$, and denote by $P$ the point of intersection of segments $AD_2$ and $BE_2$. Circle $\omega$ intersects segment $AD_2$ at two points, the closer of which to the vertex $A$ is denoted by $Q$. Prove that $AQ=D_2P$.
\end{problem}

\barydiagram{USAMO 2001/2}{2001usamo2}

\begin{soln}
  (Evan Chen) So we use barycentric coordinates.

  It's obvious that un-normalized, $D_1=(0:s-c:s-b) \implies D_2 = (0:s-b:s-c)$, so we get a normalized $D_2 = \left(0, \frac{s-b}{a}, \frac{s-c}{a}\right)$.  Similarly, $E_2 = \left(\frac{s-a}{b},0,\frac{s-c}{b}\right)$.

  Now we obtain the points $P=\left(\frac{s-a}{s}, \frac{s-b}{s}, \frac{s-c}{s}\right)$ by intersecting the lines $AD_2: (s-c)y=(s-b)z$ and $BE_2: (s-c)x=(s-a)z$.

  Let $Q'$ be such that $AQ' = PD_2$.  It's obvious that $Q'_y + P_y = A_y + D_{2y}$, so we find that $Q'_y = \frac{s-b}{a} - \frac{s-b}{s} = \frac{(s-a)(s-b)}{sa}$.  Also, since it lies on the line $AD_2$, we get that $Q'_z = \frac{s-c}{s-b} \cdot Q'_y = \frac{(s-a)(s-c)}{sa}$.  Hence, \[ Q'_x = 1 - \frac{((s-b)+(s-c))(s-a)}{sa} = \frac{sa - a(s-a)}{sa} = \frac{a}{s} \] Hence,
  \[ Q' = \left( \frac{a}{s}, \frac{(s-a)(s-b)}{sa}, \frac{(s-a)(s-c)}{sa}\right)\]

  Let $I = \left(\frac{a}{2s}, \frac{b}{2s}, \frac{c}{2s}\right)$.  We claim that, in fact, $I$ is the midpoint of $Q'D_1$.  Indeed,
  \begin{align*}
    \half \left( 0 + \frac{a}{s} \right) &= \frac{a}{2s} \\
    \half \left( \frac{(s-a)(s-b)}{sa} + \frac{s-c}{a} \right) &= \frac{(s-a)(s-b) + s(s-c) }{2sa} \\
    &= \frac{ab}{2sa} \\
    &= \frac{b}{2s} \\
    \half \left( \frac{(s-a)(s-c)}{sa} + \frac{s-b}{a} \right) &= \frac{c}{2s}
  \end{align*}

  Implying that $Q$ lies on the circle; in particular, diametrically opposite from $D_1$, so it is the closer of the two points.  Hence, $Q=Q'$, so we're done.

  %Then the displacement vector of $\overrightarrow{Q'I}$ is \[ v=Q'-I=\left(\frac{a}{2s}, \frac{2(s-a)(s-b)-ab}{2as}, \frac{2(s-a)(s-c)-ac}{2acs}\right) \] which we want to show has length $r$.  Let $j=2(s-a)(s-b)-ab$ and $k=2(s-a)(s-c)-ac$.  Using the distance formula, as well as $rs=K$, this becomes equivalent to:  \begin{align*} -\frac{1}{4} \cdot \frac{16K^2}{4s^2} &= a^2 \cdot \frac{jk}{(2as)^2} + b^2 \cdot \frac{ak}{4as^2} + c^2 \cdot \frac{aj}{4as^2} \\ &= \frac{jk + b^2k + c^2j}{4s^2} \\ \Leftrightarrow -16K^2 &= 4(jk+b^2k+c^2j) \end{align*}

  %But we also know that $-16K^2 = a^4+b^4+c^4-2a^2b^2-2b^2c^2-2c^2a^2$ by Heron, and  \begin{align*} j &= 2(s-a)(s-b)-ab \\ &= 2s^2-2(a+b)s+ab \\ &= \frac{1}{2} (a+b+c)^2 - (a+b)(a+b+c) + ab \\ &= \frac{1}{2}(a^2+b^2+c^2) + (ab+bc+ca)-(a^2+2ab+b^2+ca+ab)+ab \\ &= \frac{1}{2}(ac^2-a^2-b^2) \end{align*}  Similarly, \[ k=\frac{1}{2} (b^2-a^2-c^2)\]

  %So now we just compute  \begin{align*} 4jk + 4b^2k+4c^2j &= (c^2-a^2-b^2)(b^2-a^2-c^2) + 2b^2(b^2-a^2-c^2) + 2c^2(c^2-a^2-b^2) \\ &= (a^2+(b^2-c^2))(a^2+(c^2-b^2)) + 2b^4 + 2c^4 -2a^2b^2 - 2a^2c^2 - 4b^2c^2 \\ &= a^4-(b^2-c^2)^2  + 2b^4 + 2c^4 -2a^2b^2 - 2a^2c^2 - 4b^2c^2 \\ &= a^4+b^4+c^4-2a^2b^2 - 2b^2c^2-2c^2a^2 \\ &= -16K^2 \end{align*}  as desired.
\end{soln}

\seccm
The unusual points $D_2$, and $E_2$ make this problem ripe for barycentric coordinates, since they can be written as ratios on the appropriate sides.  The point $Q'$ is also easy to construct because of lengths.

Now it turns out that it's well known that $AD_2$ passes through the point directly above $D_1$, so it is merely a matter of length chasing to construct $Q'$ and subsequently use midpoints.

\section{USAMO 2008/2}
\label{exam:2008usamo2}
We follow this with another problem, made possible by EFFT (theorem \ref{tech:efft}).  As we shall see, EFFT is particularly useful for constructing perpendiculars to the sides of the triangle.

\secps
\begin{problem}[USAMO 2008/2]
  Let $ABC$ be an acute, scalene triangle, and let $M$, $N$, and $P$ be the midpoints of $\overline{BC}$, $\overline{CA}$, and $\overline{AB}$, respectively. Let the perpendicular bisectors of $\overline{AB}$ and $\overline{AC}$ intersect ray $AM$ in points $D$ and $E$ respectively, and let lines $BD$ and $CE$ intersect in point $F$, inside of triangle $ABC$. Prove that points $A$, $N$, $F$, and $P$ all lie on one circle.
\end{problem}

\barydiagram{USAMO 2008/2}{2008usamo2}

\begin{soln}
  (Evan Chen) Set $A=(1,0,0)$, $B=(0,1,0)$ and $C=(0,0,1)$.  Evidently $P=(\frac12, \frac12, 0)$, etc.  Check that the equation of the line $AM$ is $y=z$.

  We will now compute the coordinates of the point $D$.  Write $D=(1-2t, t, t)$ as it lies on the line $AM$.  Applying EFFT to $DP \perp AB$,
  \begin{align*}
  \left(\left(t-\frac12\right) - \left(\frac12 - 2t\right)\right)\left(-c^2/2\right) + t\left(a^2-b^2/2\right) &= 0\\
  \implies (3t-1)(-c^2) + t(a^2-b^2) &= 0 \\
  \implies (a^2-b^2-3c^2)t &= -c^2 \\
  \implies t &= \frac{c^2}{3c^2+b^2-a^2}
  \end{align*}
  Call this value $j$.  Then $D=(1-2j, j, j)$.  Similarly, $E=(1-2k, k, k)$ where $k= \frac{b^2}{3b^2+c^2-a^2}$.

  Now the line $BD$ has equation $\frac zx = \frac{j}{1-2j}$, whilst the line $CE$ has equation $\frac yx = \frac{k}{1-2k}$.  So if $F=(p,q,r)$, then $p+q+r=1$, $\frac rp = \frac{j}{1-2j}$, and $\frac qp = \frac{k}{1-2k}$.

  In fact, $\frac rp = \frac{c^2}{(3c^2+b^2-a^2)-2c^2} = \frac{c^2}{c^2+b^2-a^2}$, and $\frac qp = \frac{b^2}{b^2+c^2-a^2}$.  Evidently $\frac 1p = 1 + \frac rp + \frac qp = 1+\frac{b^2+c^2}{c^2+b^2-a^2} = \frac{2S+a^2}{S} = 2 +\frac{a^2}{S}$.  Let $S = c^2+b^2-a^2$ for convenience.

  Dilate $F$ to $F'=2F-A = (2p-1, 2q, 2r)$.  It suffices to show this lies on the circumcircle of triangle $ABC$ (since this will send $P$ and $N$ to $B$ and $C$, respectively), which has equation $a^2yz + b^2zx + c^2xy=0$.  Hence it suffices to show that
  \begin{align*}
  0 &= a^2(2q)(2r) + b^2(2r)(2p-1) + c^2(2p-1)(2q) \\
  \iff 0 &= 4a^2\frac qp \frac rp + (2+\frac1p)(2b^2 \frac rp + 2c^2 \frac qp) \\
  \iff -(2-(2+\frac{a^2}{S}))(2b^2 \frac{c^2}{S} + 2c^2 \frac{b^2}{S}) &= 4a^2 \frac{b^2c^2}{S^2}  \\
  \iff a^2(2b^2c^2 + 2c^2b^2) &= 4a^2b^2c^2 \\
  \iff 4a^2b^2c^2 &= 4a^2b^2c^2
  \end{align*}
  which is true.
\end{soln}

\seccm
EFFT provides a method for constructing the perpendicular bisectors (in actuality this is just corollary \ref{tech:perpbis}), from which it is not hard at all to intersect lines.  The final conclusion involves a circle, so we dilate the points such that this circle is simply the circumcircle, which has a very simple form.

\section{ISL 2001 G1}
\label{exam:2001islg1}
We now present an application of Conway's Formula to an ISL G1.

\secps


\begin{problem}[ISL 2001/G1]
  Let $A_1$ be the center of the square inscribed in acute triangle $ABC$ with two vertices of the square on side $BC$.  Thus one of the two remaining vertices of the square is on side $AB$ and the other is on $AC$. Points $B_1,\ C_1$ are defined in a similar way for inscribed squares with two vertices on sides $AC$ and $AB$, respectively. Prove that lines $AA_1,\ BB_1,\ CC_1$ are concurrent.
\end{problem}

\barydiagram{ISL 2001 G1}{2001islg1}

\begin{soln}
  (Max Schindler) There is an obvious homothety centered on $A$ that maps the inscribed square with center $A_1$ to one with center $A_2$, which is constructed externally off side $BC$. $B_2$ and $C_2$ can be defined similarly. It then suffices to show that $AA_2,\ BB_2,\ CC_2$ concur.

By Conway's formula, $A_2=(-a^2:S_C+S\cot{45}:S_B+S\cot{45})=(-a^2:S_C+S:S_B+S)$. It can be very easily verified that $AA_2$, and similarly $BB_2$ and $CC_2$, all go through the point $((S_B+S)(S_C+S):(S_A+S)(S_C+S):(S_B+S)(S_A+S))$ and are thus concurrent (this point is called the Outer Vecten Point).
\end{soln}
\seccm
Homothety (similar to that in problem \ref{exam:2008usamo2}) allowed us to create a convenient expression for $S$ by Conway's Formula, at which point Conway's Formula cleared the problem.

\section{2012 WOOT PO4/7}
\label{exam:2012wootpo47}
We now exhibit yet another problem in which perpendicular bisectors can be constructed via EFFT, which are then used to determine a circumcenter.  The method was used by v\_Enhance during the actual test and received $7/.7$.  See \cite{woot}.

\secps
\begin{problem}[2012 WOOT PO4/7]
  Let $\omega$ be the circumcircle of acute triangle $ABC$.  The tangents to $\omega$ at $B$ and $C$ intersect at $P$, and $AP$ and $BC$ intersect at $D$.  Points $E$, $F$ are on $AC$ and $AB$, respectively, such that $DE \parallel BA$ and $DF \parallel CA$.
  \be[(a)]
  \ii Prove that points $F$, $B$, $C$, and $E$ are concyclic.
  \ii Let $A_1$ denote the circumcenter of cyclic quadrilateral $FBCE$.  Points $B_1$ and $C_1$ are defined similarly.  Prove that $AA_1$, $BB_1$, and $CC_1$ are concurrent.
  \ee
\end{problem}

\barydiagram{WOOT Practice Olympiad 4, Problem 7}{2012wootpo47}

\begin{soln}
  (Evan Chen) Let $A=(1,0,0)$, etc.  Since the symmedian point has $K = (a^2:b^2:c^2)$, we find that $D = \left(0, \frac{b^2}{b^2+c^2}, \frac{c^2}{b^2+c^2} \right)$.  Then, it's obvious that $E = \left( \frac{b^2}{b^2+c^2}, 0, \frac{c^2}{b^2+c^2} \right)$ and $F = \left( \frac{c^2}{b^2+c^2}, \frac{b^2}{b^2+c^2}, 0 \right)$.  In that case, notice that since the general form of a circle is $a^2yz + b^2zx + c^2xy + (ux + vy + wz)(x+y+z) = 0$, the circle passing through \[ \omega: a^2yz + b^2zx + c^2xy - \frac{b^2c^2}{b^2+c^2} \cdot x(x+y+z) =0 \] passes through BFEC, so it's cyclic.

Suppose $A_1 = (x_0, y_0, z_0)$.  Let $M_1 = \frac{B+F}{2} = \left( \frac{1}{2} \frac{c^2}{b^2+c^2}, \frac{1}{2} + \frac12 \frac{b^2}{b^2+c^2}, 0 \right)$.  Then $A_1M_1 \perp AB$, applying EFFT and rearranging yields \[ x_0 - y_0 + \frac{a^2-b^2}{c^2} z_0 = \frac{c^2}{b^2+c^2} \]
Doing a similar thing, we find that $A_1$ is given by
\begin{align*}
x_0 + y_0 + z_0 &= 1 \\
x_0 - y_0 + \frac{a^2-b^2}{c^2} z_0 &= \frac{c^2}{b^2+c^2} \\
x_0 + \frac{a^2-c^2}{b^2} y_0 - z_0 &= \frac{b^2}{b^2+c^2}
\end{align*}

Then by Cramer's rule, we can bash and get $\frac{y_0}{z_0} = \frac{a^2+b^2-c^2}{a^2-b^2+c^2}$.  Then it's easy to get the desired conclusion by Ceva.
\end{soln}

\seccm
Once we realize that $AD$ is a symmedian, the symmedian point immediately gives us a nice form for $D$, and hence $E$ and $F$ by similar triangles.  The condition that $FBCE$ is cyclic is made easier since two of the points are vertices of the triangle, making both $v$ and $w$ vanish; hence the end computation is no longer hard.

For the second part, EFFT allows us to use perpendicular bisectors to construct the circumcenter; we can then compute the ratio that the cevian $AA_1$ divides $BC$ into.  Ceva's Theorem guarantees that these will cancel cyclically, so we just compute the proper determinants, knowing that they will cancel cyclically at the end.

The part about Ceva is worth remembering: any time you are trying to prove three cevians are concurrent, barycentric coordinates may provide an easy way to compute the ratios in Ceva's Theorem.

\section{ISL 2005 G5}
\label{exam:2005islg5}
We now present a solution to an ISL problem involving Strong EFFT (theorem \ref{tech:strongefft}).  This problem was proposed by Romania, and is relatively difficult to approach synthetically.

\secps
\begin{problem}[ISL 2005/G5]
  Let $\triangle ABC$ be an acute-angled triangle with $AB \neq AC$. Let $H$ be the orthocenter of triangle $ABC$, and let $M$ be the midpoint of the side $BC$. Let $D$ be a point on the side $AB$ and $E$ a point on the side $AC$ such that $AE=AD$ and the points $D$, $H$, $E$ are on the same line. Prove that the line $HM$ is perpendicular to the common chord of the circumscribed circles of triangle $\triangle ABC$ and triangle $\triangle ADE$.
\end{problem}

\barydiagram{ISL 2005 G5}{2005islg5}

\begin{soln}
  (Max Schindler) We use barycentric coordinates.

  Let the length $AD=AE=\ell$.

  Then, $D=(c-\ell:\ell:0)$, $E=(b-\ell:0:\ell)$, and, using Conway's Notation, $H=\left(\frac{1}{S_A}:\frac{1}{S_B}:\frac{1}{S_C}\right)$.

  Since they are collinear,
  \begin{align*}
  \det \left(
  \begin{array}{ccc}
  c-\ell & \ell & 0 \\
  b-\ell & 0 & \ell \\
  \frac{1}{S_A} & \frac{1}{S_B} & \frac{1}{S_C}
  \end{array}
  \right) &=0 \\
  \implies \frac{-\ell(c-\ell)}{S_B}-\ell\left(\frac{b-\ell}{S_C}-\frac{\ell}{S_A}\right) &= 0 \\
  \implies \frac{cS_C-bS_B}{S_BS_C} &= \ell\left[\frac{1}{S_A}+\frac{1}{S_B}+\frac{1}{S_C}\right]
  \end{align*}



  Thus, $\ell = \frac{S_A(cS_C-bS_B)}{S_BS_C+S_AS_C+S_AS_B}$.

  Substituting in $S_A=\frac{b^2+c^2-a^2}{2}$ and similar, remembering that $S_BS_C+S_AS_C+S_AS_B=S^2$ gives:

  \[ \ell=\frac{(b^2+c^2-a^2)[c(a^2+b^2-c^2)+b(a^2-b^2+c^2)]}{(a+b+c)(a+b-c)(b+c-a)(c+b-a)}\]

  It can easily be seen that $c=-b$ is a root of the bracketed part of the numerator; factoring it out leaves $a^2+2bc-b^2-c^2=a^2-(b-c)^2=(a-b+c)(a+b-c)$.

  Thus, \[ \ell=\frac{(b^2+c^2-a^2)(b+c)(a-b+c)(a+b-c)}{(a+b+c)(a+b-c)(b+c-a)(c+b-a)} \implies \boxed{\ell = \frac{(b^2+c^2-a^2)(b+c)}{(a+b+c)(b+c-a)}}\]




  Meanwhile, the circumcircle of $ADE$ has equation given by $-a^2yz-b^2zx-c^2xy + (x+y+z)(ux+vy+wz) = 0$, as usual, so upon substituting the points $A$, $D$ and $E$ we find that

  \begin{align*}
     u &= 0 \\
     -c^2(c-\ell)(\ell) + c\left((c-\ell) u + \ell v\right) &= 0 \\
     -b^2(b-\ell)(\ell) + b\left((b-\ell) u + \ell w\right) &= 0
  \end{align*}

  Conveniently enough, the first equation immediately gives $u=0$, and we can easily solve for $v$ and $w$ to get $v = c(c-\ell)$ and $w=b(b-\ell)$.

  Now we need the common chord of the two circles

  \begin{align*}
     -a^2yz - b^2zx - c^2xy &= 0 \\
     -a^2yz - b^2zx - c^2xy + (x+y+z)(c(c-\ell) y + b(b-\ell) z) &= 0\\
  \end{align*}

  But this is easy: subtract the two equations.  We find that the common chord has equation
  \[ c(c-\ell) y + b(b-\ell) z = 0\]




  Now, we have everything we need to finish this problem off.

  Consider two points on this chord, $A=(1,0,0)$ and $P=(0:b(b-\ell):-c(c-\ell))$.

  Then, the (scaled) displacement vector $\overrightarrow{PA}$ is $(b(b-\ell)-c(c-\ell):-b(b-\ell):c(c-\ell))$.

  Let the circumcenter be the null vector. Then, the displacement vector $\overrightarrow{MH}$ is $\vec A+\vec B + \vec C-\frac{\vec B + \vec C}{2}=\left(1,\frac{1}{2},\frac{1}{2}\right)$ which can be scaled as $(2:1:1)$

  Since $PA$ has coefficient-sum $0$, we can apply Strong EFFT: \[ MH \perp PA \Leftrightarrow a^2[c(c-\ell)-b(b-\ell)]+b^2[c(c-\ell)+b(b-\ell)]+c^2[-b(b-\ell)-c(c-\ell)] =0 \] which becomes  \[ b(a^2-b^2+c^2)(b-\ell)=c(a^2+b^2-c^2)(c-\ell) \]

  But we have \[ b-\ell=\frac{b(a+b+c)(b+c-a)-(b^2+c^2-a^2)(b+c)}{(a+b+c)(b+c-a)}=\frac{c(a^2+b^2-c^2)}{(a+b+c)(b+c-a)}\]  Similarly, \[ c-\ell=\frac{b(a^2-b^2+c^2)}{(a+b+c)(b+c-a)} \]

  Thus,
  \begin{align*}
    b(a^2-b^2+c^2)(b-\ell) &= b(a^2-b^2+c^2)\frac{c(a^2+b^2-c^2)}{(a+b+c)(b+c-a)} \\
    &=c(a^2+b^2-c^2)(c-\ell) \\
    &\implies MH \perp PA
  \end{align*}
\end{soln}

\seccm
While we eventually did have to use the somewhat messy form of $H$, Strong EFFT permits us to avoid doing so when finding the condition for perpendicular chords.  The other useful insight was computing the radical axis by simply subtracting the two equations: keep this in mind!

By the way, it is true in general (as the diagram implies) that $M$, $H$ and $F$ are collinear.  To my best knowledge, the easiest way to prove this is to first prove the original problem.

\section{USA TSTST 2011/4}
\label{exam:2011tstst4}
The solution to this problem involves a synthetic observation, followed by an application of Vieta.   Normally, solving for $P$ involves an incredibly painful quadratic.  However, by first computing the other intersection, we can actually use Vieta to deduce the coordinate of the point we care about.
\secps
\begin{problem}[TSTST Problem 4]
  Acute triangle $ABC$ is inscribed in circle $\omega$. Let $H$ and $O$ denote its orthocenter and circumcenter, respectively. Let $M$ and $N$ be the midpoints of sides $AB$ and $AC$, respectively. Rays $MH$ and $NH$ meet $\omega$ at $P$ and $Q$, respectively. Lines $MN$ and $PQ$ meet at $R$. Prove that $OA \perp RA$.
\end{problem}

\barydiagram{2011 TSTST, Problem 4}{2011tstst4}

\begin{soln}
  (Evan Chen) We use barycentric coordinates.  Set $A=(1,0,0)$, $B=(0,1,0)$ and $C=(0,0,1)$.  Let $T_A = \frac{1}{S_A}$ in Conway's Notation, so that $H = (T_A : T_B : T_C)$.

  First we will compute the point $R' = (x:y:z)$ on $MN$ with $R'A \perp OA$.  Setting $\vec O = 0$, and applying Strong EFFT, we find that the point satisfies equations $c^2y+b^2z=0$ (per EFFT) and $x=y+z$ (per midline).  Thus, we may write $\boxed{R' = (b^2-c^2:b^2:-c^2)}$.

  Next, we will compute the coordinates of $P$.  Check that we have the line \[ HM: x-y + \frac{T_B - T_A}{T_C} z = 0 \] Furthermore, the circumcircle has equation $z[a^2y+b^2x] + c^2xy = 0$.  WLOG, homogenize so that $z=-T_C$; then the first equation becomes $y = x + T_A - T_B$.  Substituting in the second equation, this becomes simply
  \[ c^2x(x+(T_A-T_B)) + a^2(-T_C)(x+T_A-T_B) + b^2(-T_C)(x) = 0 \]
  Expanding, and collecting coefficients of $x$, this reduces to
  \[ c^2x^2 + [c^2(T_A-T_B)-a^2T_C-b^2T_C] x + a^2(-T_C)(T_A-T_B) = 0 \]

  Now here's the trick: it's well known that the reflection of $H$ across $M$ lies on the circle.  So the other solution to this equation is $\vec H' = \vec A + \vec B - \vec H$ (since $H'AHB$ is a parallelogram).  We can thus quickly deduce that, in barycentric coordinates, $H'=(T_B+T_C : T_A+T_C : -T_C)$.  So by Vieta's Formulas, the $x$-coordinate we seek is just
  \[ x = -\frac{c^2(T_A-T_B)-a^2T_C-b^2T_C}{c^2} - (T_B+T_C) = \frac{a^2+b^2-c^2}{c^2} T_C - T_A \]
  Hence, \[ y = \frac{a^2+b^2-c^2}{c^2} T_C - T_B \]

  So finally, we obtain \[ \boxed{P = \left( \frac{a^2+b^2-c^2}{c^2} T_C - T_A : \frac{a^2+b^2-c^2}{c^2}T_C - T_B : -T_C \right)} \] Similarly, \[ \boxed{Q = \left( \frac{a^2+c^2-b^2}{b^2} T_B - T_A : -T_B : \frac{a^2+c^2-b^2}{b^2}T_B - T_C \right) } \]

  It remains to show that $P$, $Q$ and $R'$ are collinear.  This occurs if and only if
  \[ 0 = \det \left( \begin{array}{ccc}
  \frac{a^2+b^2-c^2}{c^2} T_C - T_A & \frac{a^2+b^2-c^2}{c^2}T_C - T_B & -T_C \\
  \frac{a^2+c^2-b^2}{c^2} T_B - T_A & -T_B & \frac{a^2+c^2-b^2}{b^2}T_B - T_C  \\
  b^2-c^2 & b^2 & -c^2 \end{array} \right) \]

  Subtracting the second and third column from the first, this is equivalent to
  \[ 0 = \det \left( \begin{array}{ccc}
  T_C + T_B - T_A & \frac{a^2+b^2-c^2}{c^2}T_C - T_B & -T_C \\
  T_C + T_B - T_A & -T_B & \frac{a^2+c^2-b^2}{b^2}T_B - T_C  \\
  0 & b^2 & -c^2 \end{array} \right) \]

  Expanding the determinant, and factoring out the term $T_C+T_B-T_A$, we see that it suffices to show that
  \[ 0 = c^2T_B - b^2T_C - \left[ -(c^2)\left(\frac{a^2+b^2-c^2}{c^2} T_C - T_B\right) + (b^2)\left(\frac{a^2+c^2-b^2}{b^2} T_B - T_C \right) \right]\]
  Expanding, this is just
  \[ -(a^2-b^2+c^2) T_B + (a^2+b^2-c^2) T_C \]

  But $T_B = \frac{1}{S_B} = \frac{2}{a^2-b^2+c^2}$ and $T_C = \frac{2}{a^2+b^2-c^2}$.  Since $-2 + 2 = 0$, we conclude that the three points are collinear.

\seccm
The application of Vieta played a critical role in this problem.  BTW, during the actual TSTST, only one person successfully bashed this problem (via complex numbers).  Had I known about this technique at the time, I think I could have gotten a positive score on Day 2.

\section{USA TSTST 2011/2}
Finally, to complete the section, we present an example of a bash which is possibly not feasible by hand.  Incidentally, this is the only problem in the section that was not done by hand.  Max Schindler claims this is doable if you have the computational fortitude.

The problem has been inverted around the point $A$, and the constant $c$ has been replaced by $1/c$.

\secps
\begin{problem}[TSTST Problem 2, Inverted]
  Let $APQ$ be an obtuse triangle with $\angle A > 90\dg$.  The tangents at $P$ and $Q$ to the circumcircle of $APQ$ meet at the point $B$.  Points $X$ and $Y$ are selected on segments $PB$ and $QS$, respectively, such that $PX = c \cdot QA$, and $QY = c \cdot PA$.  Then, $PA$ and $QA$ are extended through $A$ to $R$ and $S$, so that $AR = AS = c \cdot \frac{PA \cdot QA}{PQ}$.  $O$ is the point such that $AROS$ is a parallelogram.   Given that $AXOY$ is cyclic, prove that $\angle PXA = \angle QYA$.
\end{problem}

\barydiagram{2011 TSTST, Problem 2}{2011tstst2}

\begin{soln}
  The reference triangle is $APQ$, with $A=(1,0,0)$, $P=(0,1,0)$ and $Q=(0,0,1)$.  To avoid confusion, let $a=PQ$, $q=AP$ and $p=AQ$.

  It is not hard at all to compute $R$ and $S$.  Since $AR = AS = \frac{cpq}{a}$, and $A=(1,0,0)$, $P=(0,1,0)$ and $Q=(0,0,1)$, a quick application of ratios of lengths yields \[ R=\left(\frac{c q}{a}+1,\frac{-c q}{a},0\right) \] and \[ S=\left(\frac{c p}{a}+1,0,\frac{-c p}{a}\right)\]

  Next we want to compute $X=(x,y,z)$.  Since $X$ lies on the tangent $PB$, it satisfies $0=b^2y+c^2z$ (see the earlier problem, \ref{exam:2011tstst4}).  Letting $x+y+z=1$, and invoking the distance formula, we find that $(cq)^2 = a^2(y-1)z + b^2zx + c^2x(y-1)$.  Solving this system by writing everything in $z$, and taking the negative root, gives $X = \left( \frac{a c}{p},\frac{-a^2 c+a p+c q^2}{a p},-\frac{c q^2}{a p} \right)$.

  Similarly, $Y = \left( \frac{a c}{q},-\frac{c p^2}{a q},\frac{a^2 (-c)+a q+c p^2}{a q} \right)$.  We then obtain that $O = \left( 1+\frac{c p}{a}+\frac{c q}{a},-\frac{c p}{a},-\frac{c q}{a}\right)$

  We'd like to find criteria for when $\angle PXA = \angle QYA$.  Noting the equal angles $\angle BPQ = \angle BQP$, we extend $XA$ and $YA$ through $A$ to points $X'$ and $Y'$ respectively.  Then $X' = (0 : pa + q^2c - a^2c : -q^2c)$ and $Y'=(0 : -p^2c : qa + p^2c - a^2 c)$ (aren't cevians wonderful?).  In that case, we obtain the lengths \[ PX' = \frac{-q^2c}{pa-a^2c} a = \frac{-q^2c}{p-ac} \] and similarly, \[ QY' = \frac{-p^2c}{q-ac} \]

  Now $\angle PXA = \angle QYA$ if and only if the triangles $XPX'$ and $YQY'$ are similar!  Recall that $PX=cq$ and $QY=cp$; by SAS, this occurs if and only if
  \[ \frac{ \frac{-q^2c}{p-ac} }{cq} = \frac{ \frac{-p^2c}{q-ac} }{cp} \iff q(q-ac) = p(p-ac) \iff c = \frac{p+q}{a}\]

  So we just need to show that if $AROS$ is cyclic, then $c=\frac{p+q}{a}$.

  \bigskip \bigskip

  Unfortunately, this is where the calculations get heavy.\footnote{I wonder if there's any synthetic interpretation that would lead directly to this expression for $c$.  That would be nice too.  Tell me if you find one!}

  Putting $A,X,Y$ into the circle formula and solving for $u$, $v$, $w$ yields
  \begin{align*}
    u &= 0 \\
    v &= (c^2 p q^2 (-a^2 c + a q + c p (p + q)))/(-a^3 c^2 + a^2 c (p + q) - c (p^3 + q^3) + a (-p q + c^2 (p^2 + q^2))) \\
      w &= ( c^2 p^2 q (-a^2 c + a p + c q (p + q)))/(-a^3 c^2 + a^2 c (p + q) - c (p^3 + q^3) + a (-p q + c^2 (p^2 + q^2)))
  \end{align*}

  Substituting the point $O$ into the equation and simplifying, we get that
  \[ \scriptstyle { \frac{-(c p q (a^5 c^3 - 2 a^4 c^2 (p + q) -
       c^2 (p + q)^3 (p^2 - p q + q^2) +
       a^3 c ((1 - 2 c^2) p^2 + 3 p q + (1 - 2 c^2) q^2) +
       a c (p + q)^2 ((-1 + c^2) p^2 - c^2 p q + (-1 + c^2) q^2) +
       a^2 (p + q) (-p q + 3 c^2 (p^2 + q^2))))}{
       (a^2 (a^3 c^2 -
       a^2 c (p + q) + c (p^3 + q^3) - a (-p q + c^2 (p^2 + q^2))))} = 0 }\]

    Clearing the denominator, this is just
    \begin{align*}
  0 =& -c p q (a^5 c^3 - 2 a^4 c^2 (p + q)  \\
   & -c^2 (p + q)^3 (p^2 - p q + q^2) \\
   & + a^3 c ((1 - 2 c^2) p^2 + 3 p q + (1 - 2 c^2) q^2) \\
   & + a c (p + q)^2 ((-1 + c^2) p^2 - c^2 p q + (-1 + c^2) q^2) \\
   & + a^2 (p + q) (-p q + 3 c^2 (p^2 + q^2)))
  \end{align*}

   Noting that $c=0$ and $c=\frac{p+q}{a}$ are solutions (the latter ``by inspection'') we can divide this out and collect terms in $c$ to get
   \[ c^2 \left(-a^4-p^4-p^3 q-p q^3-q^4+2 a^2 \left(p^2+q^2\right)\right)+c \left(a^3 (p+q)-a \left(p^3+p^2 q+p q^2+q^3\right)\right) - a^2pq = 0 \]

   We need to show that there are no positive reals $c$ satisfying this.  A quick glance at the constant term tells us that we should probably use the discriminant; our only hope is that there are no real roots at all.  Computing, we get
   \begin{align*}
   \Delta =& a^6 p^2 - 2 a^4 p^4 + a^2 p^6 - 2 a^6 p q + 4 a^4 p^3 q - 2 a^2 p^5 q \\ &+ a^6 q^2 - 4 a^4 p^2 q^2 - a^2 p^4 q^2 + 4 a^4 p q^3 + 4 a^2 p^3 q^3 \\ &- 2 a^4 q^4 - a^2 p^2 q^4 - 2 a^2 p q^5 + a^2 q^6 \end{align*}
   which factors as
   \[ \Delta = -a^2 (-a+p+q) (p-q)^2 (a+p-q) (a-p+q) (a+p+q) \]

   Since $a$, $p$ and $q$ are the sides of a triangle, this is always negative.

\end{soln}

\seccm
The calculations are relatively easy up until incorporating the condition that $AROS$ is cyclic.  Despite the fact that the ``cyclic'' condition was fairly brutal, it illustrates a few points worth making:
\begin{parlist}
  \item similar triangles are your friends in barycentrics too, because they can yield equal angles,
  \item inversion can provide a way to eliminate awkward circle conditions,
  \item unusual side length conditions related to the sides of the triangles is not a problem, and
  \item cyclic quadrilaterals suck.
\end{parlist}

Please let us know if you find a cyclic criteria which completes a reasonable solution!  We'd like to have a human to attribute the solution to...




\end{soln}

\chapter{Problems}
\begin{quote}
  There is light at the end of the tunnel, but it is moving away at speed $c$.
\end{quote}


Here, we provide some problems for the reader to attempt.  Some are easier than others.  Infinity \cite{infinity} has a longer list than is included here; see \url{http://www.bmoc.maths.org/home/areals.pdf} for the relevant excerpt.
%\newcommand{\sollink}[1]{\footnote{A rough outline of the solution can be found on the AoPS thread \href{#1}{here}.}}
\newcommand{\sollink}[1]{}

For practice ``in real life'', i.e. determining when barycentric coordinates are feasible for a problem (here, they are), you can often simply take an arbitrary geometry problem that begins with the words ``triangle $ABC$'', and see if it's possible to get a barycentric solution.  Some will, of course, take much longer than others, but a good portion will have a reasonable solution.

\section{Problems}

\begin{enumerate}
  \item Prove Stewart's Theorem.

  \item Prove Routh's Theorem.

  \item (USA TST 2003/2) Let $ABC$ be a triangle and let $P$ be a point in its interior. Lines $PA$, $PB$, $PC$ intersect sides $BC$, $CA$, $AB$ at $D$, $E$, $F$, respectively. Prove that
\[ [PAF]+[PBD]+[PCE]=\frac{1}{2}[ABC]  \]
if and only if $P$ lies on at least one of the medians of triangle $ABC$. (Here $[XYZ]$ denotes the area of triangle $XYZ$.)

  \item (Mongolia TST 2000/6) In a triangle $ABC$, the angle bisector at $A$,$B$,$C$ meet the opposite sides at $A_1$,$B_1$,$C_1$, respectively. Prove that if the quadrilateral $BA_1B_1C_1$ is cyclic, then
  \[ \frac{AC}{AB+BC} = \frac{AB}{AC+BC} + \frac{BC}{BA+AC} \]

  \item (ISL 2005/G1) Given a triangle $ABC$ satisfying $AC+BC=3\cdot AB$. The incircle of triangle $ABC$ has center $I$ and touches the sides $BC$ and $CA$ at the points $D$ and $E$, respectively. Let $K$ and $L$ be the reflections of the points $D$ and $E$ with respect to $I$. Prove that the points $A$, $B$, $K$, $L$ lie on one circle.  \sollink{http://www.artofproblemsolving.com/Forum/viewtopic.php?p=2351823\#p2351823}

  \item (APMO 2005/5) In a triangle $ABC$, points $M$ and $N$ are on sides $AB$ and $AC$, respectively, such that $MB = BC = CN$. Let $R$ and $r$ denote the circumradius and the inradius of the triangle $ABC$, respectively. Express the ratio $MN/BC$ in terms of $R$ and $r$.  \sollink{http://www.artofproblemsolving.com/Forum/viewtopic.php?f=47\&t=30955\&p=2376234\#p2376234}

  \item (ISL 2008/G4) In an acute triangle $ ABC$ segments $ BE$ and $ CF$ are altitudes. Two circles passing through the point $ A$ and $ F$ and tangent to the line $ BC$ at the points $ P$ and $ Q$ so that $ B$ lies between $ C$ and $ Q$. Prove that lines $ PE$ and $ QF$ intersect on the circumcircle of triangle $ AEF$.

  \item (ISL 2001/G6) Let $ABC$ be a triangle and $P$ an exterior point in the plane of the triangle. Suppose the lines $AP$, $BP$, $CP$ meet the sides $BC$, $CA$, $AB$ (or extensions thereof) in $D$, $E$, $F$, respectively. Suppose further that the areas of triangles $PBD$, $PCE$, $PAF$ are all equal. Prove that each of these areas is equal to the area of triangle $ABC$ itself.  \sollink{http://www.artofproblemsolving.com/Forum/viewtopic.php?p=119203\&sid=f14625806c6d8ca5231f857e4519121e\#p119203}

  \item (IMO 2007/2) Consider five points $A$, $B$, $C$, $D$ and $E$ such that $ABCD$ is a parallelogram and $BCED$ is a cyclic quadrilateral. Let $\ell$ be a line passing through $A$. Suppose that $\ell$ intersects the interior of the segment $DC$ at $F$ and intersects line $BC$ at $G$. Suppose also that $EF = EG = EC$. Prove that $\ell$ is the bisector of angle $DAB$. \sollink{http://www.artofproblemsolving.com/Forum/viewtopic.php?f=451\&t=159837\&p=2660858\#p2660858}
\end{enumerate}


\section{Hints}
Hints have been obfuscated by ROT13.
\begin{enumerate}
  \item ABG UNEQ
  \item FGENVTUGSBEJNEQ
  \item HFR NERN SBEZ
  \item PVEPYR SBEZ FBYIR HIJ
  \item PVEPF GURA J VF ZVAHF N O
  \item QVFG SNPGBE GURA TRB VQ
  \item CBJRE CBVAG FUBJF NRCD PLPYVP GURA NATYR PUNFR
  \item PNERSHY FVTARQ NERNF
  \item ERS OPQ... SVAQ S T FLF RDA RYVZ P M SNPGBE
\end{enumerate}

\chapter{Additional Problems from MOP 2012}
\begin{quote}
  If our solution has at table of contents, how should we number the pages?
  \\ -- Evan Chen, MOP 2012
\end{quote}

Here's a more recent collection of geometry problems solvable using barycentric coordinates.  These are all problems which appeared at MOP 2012,\footnote{MOP, occasionally called MOSP, stands for Mathematical Olympiad Program; it is the USA training program for the IMO.} along with an IMO question from that year.

These are in very roughly ascending order of difficulty, based both on the amount of computation and the difficulty of setup.
\begin{enumerate}
  \ii (IMO 2012) Given triangle $ABC$ the point $J$ is the center of the excircle opposite the vertex $A$. This excircle is tangent to the side $BC$ at $M$, and to the lines $AB$ and $AC$ at $K$ and $L$, respectively. The lines $LM$ and $BJ$ meet at $F$, and the lines $KM$ and $CJ$ meet at $G.$ Let $S$ be the point of intersection of the lines $AF$ and $BC$, and let $T$ be the point of intersection of the lines $AG$ and $BC.$ Prove that $M$ is the midpoint of $ST.$
  \ii (USA TST 2012) In acute triangle $ABC$, $\angle A < \angle B$ and $\angle A < \angle C$.  Let $P$ be a variable point on side $BC$.  Points $D$ and $E$ lie on sides $AB$ and $AC$, respectively, such that $BP = PD$ and $CP = PE$.  Prove that if $P$ moves along side $BC$, the circumcircle of triangle $ADE$ passes through a fixed point other than $A$.
  \ii (ELMO\footnote{This year's test was called ``Every Little Mistake $\implies$ 0''.} 2012/5) Let $ABC$ be an acute triangle with $AB < AC$, and let $D$ and $E$ be points on side $BC$ such that $BD$ = $CE$ and $D$ lies between $B$ and $E$.  Suppose there exists a point $P$ inside $ABC$ such that $PD \parallel AE$ and $\angle PAB = \angle EAC$.  Prove that $\angle PBA = \angle PCA$.
  \ii (Saudi Arabia TST 2012) Let $ABCD$ be a convex quadrilateral such that $AB = AC = BD$.  The lines $AC$ and $BD$ meet at point $O$, the circles $ABC$ and $ADO$ meet again at point $P$, and the lines $AP$ and $BC$ meet at point $Q$.  Show that $\angle COQ = \angle DOQ$.
  \ii (ISL 2011/G2)  Let $A_1A_2A_3A_4$ be a non-cyclic quadrilateral.  For $1 \le i \le 4$, let $O_i$ and $r_i$ be the circumcenter and the circumradius of triangle $A_{i+1}A_{i+2}A_{i+3}$ (where $A_{i+4} = A_i$).  Prove that
    \[ \frac{1}{O_1A_1^2-r_1^2} + \frac{1}{O_2A_2^2-r_2^2} + \frac{1}{O_3A_3^2-r_3^2} + \frac{1}{O_4A_4^2-r_4^2} = 0 \]
  \ii (USA TSTST 2012) Triangle $ABC$ is inscribed in circle $\Omega$.  The interior angle bisector of angle $A$ intersects side $BC$ and $\Omega$ at $D$ and $L$ (other than $A$), respectively.  Let $M$ be the midpoint of side $BC$.  The circumcircle of triangle $ADM$ intersects sides $AB$ and $AC$ again at $Q$ and $P$ (other than $A$), respectively.  Let $N$ be the midpoint of segment $PQ$, and let $H$ be the foot of the perpendicular from $L$ to line $ND$.  Prove that line $ML$ is tangent to the circumcircle of triangle $HMN$.
  \ii (ISL 2011/G4) Acute triangle $ABC$ is inscribed in circle $\Omega$.  Let $B_0$, $C_0$, $D$ lie on sides $AC$, $AB$, $BC$, respectively, such that $AB_0 = CB_0$, $AC_0 = BC_0$, and $AD \perp BC$.  Let $G$ denote the centroid of triangle $ABC$, and let $\omega$ be a circle through $B_0$ and $C_0$ that is tangent to $\Omega$ at a point $X$ other than $A$.  Prove that $D$, $G$, $X$ are collinear.
  \ii Let $ABC$ be a triangle with circumcenter $O$ and let the angle bisector of $\angle BAC$ intersect $BC$ at $D$.  The point $M$ is such that $MC \perp BC$ and $MA \perp AD$.  Lines $BM$ and $OA$ intersect at the point $P$.  Show that the circle centered at $P$ and passing through $A$ is tangent to segment $BC$.
  \ii (USA TSTST 2012) Let $ABCD$ be a quadrilateral with $AC = BD$.  Diagonals $AC$ and $BD$ meet at $P$.  Let $\omega_1$ and $O_1$ denote the circumcircle and circumcenter of triangle $ABP$.  Let $\omega_2$ and $O_2$ denote the circumcircle and circumcenter of triangle $CDP$.  Segment $BC$ meets $\omega_1$ and $\omega_2$ again at $S$ and $T$ (other than $B$ and $C$), respectively.  Let $M$ and $N$ be the midpoints of minor arcs $SP$ (not including $B$) and $TP$ (not including $C$).  Prove that $MN \parallel O_1O_2$.
\end{enumerate}


\appendix
\chapter{Vectors Tutorial}
\label{appendix:vectors}
\begin{quote}
  ``Are you making progress on these problems? Or are you making Congress?''
  \\ --Zuming Feng
\end{quote}

\section{Definitions}
Our definition of a vector is just an arrow that points from some reference point, denoted $\vec 0$, to another point.  In this paper, we are only considering vectors in the plane (e.g. two-dimensional vectors), but the concept extends easily to multiple dimensions.  It is assumed that the reader is familiar with vector addition, the concept of magnitude, etc.

In order to make the theory of vectors applicable to geometry, we can associate every point with the vector that points to it.

Of particular use in the world of vectors is the dot product defined by
\begin{definition}
  Let $\vec u$ and $\vec v$ be vectors.  Then the \textit{dot product} $\vec u \cdot \vec v$ is defined as
  \[ \vec u \cdot \vec v = |\vec u||\vec v| \cos \theta \]
  where $\theta$ is the angle between the two vectors.
\end{definition}

The dot product has the following very nice properties:
\begin{enumerate}
  \item $\vec u \cdot \vec u = |u|^2$.
  \item $\vec u \cdot \vec v = 0$ if and only if the vectors are perpendicular.
  \item The dot product is commutative, distributive, and associative.
\end{enumerate}

\begin{exercise}
  Let $A,B,C,D$ be points in the plane.  Prove that $AC \perp BD$ if and only if $(\vec A - \vec C) \cdot (\vec B - \vec D) = 0$.
\end{exercise}

\section{Triangle ABC}
Let $ABC$ be a triangle in the plane with circumcenter $O$.  Let $\vec O = \vec 0$ be the reference point.  Then it's obvious that $|\vec A| = |\vec B| = |\vec C| = R$, where $R$ is the circumradius of the triangle.

Now for the key fact:
\begin{fact}
  $\vec A \cdot \vec B = R^2 - \frac{c^2}{2}$.  Cyclic variations hold.
\end{fact}
\begin{proof}
  Using the definition,
  \begin{align*}
    \vec A \cdot \vec B &= | \vec A || \vec B | \cos \angle AOB \\
    &= R^2 \cos 2\angle ACB \\
    &= R^2 (1 - 2\sin^2 \angle ACB) \\
    &= R^2 (1 - 2\sin^2 C) \\
    &= R^2 - \half (2R \cdot \sin C)^2 \\
    &= R^2 - \frac{c^2}{2}
  \end{align*}
  the last step following by the Law of Sines.
\end{proof}

Barycentric coordinates allow us to write any point in the plane in the form $x \vec A + y \vec B + z \vec C$.  In that case, displacement vectors, points, etc. can all be written as a linear combination of $\vec A$, $\vec B$ and $\vec C$.  The above fact allows us to compute arbitrary dot products via the distributive law.  Using properties 1 and 2 above, we can construct criterion for perpendiculars and lengths; this leads to EFFT and the distance formula.

\section{Special Points}
Again setting $\vec O = \vec 0$, we have the orthocenter\footnote{To check that this is correct, verify that $(\vec H - \vec A) \cdot (\vec B - \vec C) = 0$, implying $HA \perp BC$.} $\vec H = \vec A + \vec B + \vec C$, the centroid $\vec G = \frac{1}{3} \vec H$, and the nine-point center $\vec N = \half \vec H$.  In other words, the points on the Euler Line are very, very nice.

This allows us, for example, to compute $OH^2$ as
\begin{align*}
  OH^2 &= \vec H \cdot \vec H \\
  &= (\vec A + \vec B + \vec C) \cdot (\vec A + \vec B + \vec C) \\
  &= 3R^2 + 2(R^2-a^2/2) + 2(R^2-b^2/2) + 2(R^2-c^2/2) \\
  &= 9R^2-a^2-b^2-c^2
\end{align*}


\chapter{Formula Sheet}
\label{appendix:formulasheet}
\begin{quote}
  My formula sheets were pretty thorough, but perhaps they were missing something.
  \\ --Richard Rusczyk
\end{quote}

%import tech
\section{Standard Formulas}

\makeatletter
\let\actuallabel\label
\makeatother
\renewcommand{\label}[1]{
%Do Nothing
}
\InputTempFile{tech}
\renewcommand{\label}[1]{\label{#1}}


\section{More Obscure Formulas}
Here's some miscellaneous formulas and the like.  These were not included in the main text.

From \cite{leibnizthm},
\begin{theorem}[Leibniz Theorem]
  Let $Q$ be a point with homogeneous barycentric coordinates $(u: v: w)$ with respect to $\triangle ABC$.  For any point $P$ on the plane $ABC$ the following relation holds:
  \[ uPA^2 + vPB^2 + wPC^2 = (u + v + w)PQ^2 + uQA^2 + vQB^2 + wQC^2 \]
\end{theorem}

\subsection{Other Special Points}
\begin{tabular}{r|l}
  Point & Coordinates \\ \hline
  Gregonne Point \cite{wolfram} & $\text{Ge} = ((s-b)(s-c) : (s-c)(s-a) : (s-a)(s-b))$ \\
  Nagel Point \cite{wolfram} & $\text{Na} = (s-a : s-b : s-c)$ \\
  Isogonal Conjugate \cite{infinity} & $P^\ast = \left( \frac{a^2}{x} : \frac{b^2}{y} : \frac{c^2}{z} \right)$ \\
  Isotomic Conjugate \cite{infinity} & $P^t = \left( \frac{1}{x} : \frac{1}{y} : \frac{1}{z} \right)$ \\
  Feuerbach Point \cite{fpoint} & $F = \left((b+c-a)(b-c)^2 : (c+a-b)(c-a)^2 : (a+b-c)(a-b)^2 \right)$ \\
  Nine-point Center & $N = (a\cos(B-C) : b\cos(C-A) : c\cos(A-B))$
\end{tabular}

\subsection{Special Lines and Circles}

\begin{tabular}{rl}
  Nine-point Circle & $-a^2yz-b^2xz-c^2xy+ \half (x+y+z)(S_Ax+S_By+S_Cz) = 0$ \\
  Incircle & $-a^2yz - b^2zx  - c^2xy + (x+y+z)\left( (s-a)^2x + (s-b)^2y + (s-c)^2z \right) = 0$ \\
  $A$-excircle \cite{book} & $-a^2yz - b^2zx - c^2xy + (x+y+z)\left( s^2 x + (s-c)^2 y + (s-b)^2 z \right) = 0$ \\
  Euler Line \cite{book} & $S_A(S_B-S_C)x + S_B(S_C-S_A)y + S_C(S_A-S_B)z = 0$
\end{tabular}


%%BIBLIOGRAPHY
\begin{thebibliography}{99}
\iflong
  \addcontentsline{toc}{chapter}{\bibname}
\else
  %\addcontentsline{toc}{section}{References}
\fi

\bibitem{infinity}
  Hojoo Lee, Tom Lovering, and Cosmin Pohoata.
  \textit{Infinity}. \\
  Relevant excerpt at \url{http://www.bmoc.maths.org/home/areals.pdf}.

\bibitem{abel}
  Zachary Abel.
  ``Barycentric Coordinates.'' \\
  \url{http://zacharyabel.com/papers/Barycentric\_A07.pdf}

\bibitem{wolfram}
  Eric W. Weisstein.
  ``Barycentric Coordinates.''
  From MathWorld--A Wolfram Web Resource.
  \url{http://mathworld.wolfram.com/BarycentricCoordinates.html}

\bibitem{spanish}
  Francisco J. Garc\'ia Capit\'an.
  ``Coordenadas Baric\'entricas''. \\
  \url{http://www.aloj.us.es/rbarroso/trianguloscabri/sol/sol202garcap/cb.pdf}

\bibitem{yiu}
  Paul Yiu.
  ``The use of homogeneous barycentric coordinates in plane euclidean geometry''. \\
  \url{http://math.fau.edu/Yiu/barycentricpaper.pdf}

\bibitem{woot}
  AoPS.  Practice Olympiad 4.  WOOT 2011-2012.

\bibitem{leibnizthm}
  AoPS Forum.  \\ \url{http://www.artofproblemsolving.com/Forum/viewtopic.php?p=2162335#p2162335}

\bibitem{fpoint}
  AoPS Forum.  \\ \url{http://www.artofproblemsolving.com/Forum/viewtopic.php?f=46&t=399047&p=2225645}

\bibitem{book}
  Christopher J. Bradley.  \textit{Challenges in Geometry: For Mathematical Olympians Past and Present}. \\


\end{thebibliography}



%\TempToTeX{tech}

\end{document}
